%% kapa: exam=1 round=22 year=2011 period=2 form=A subject=appraisal_law subject_ko=부동산관계법규 seq=39 number=79
%% source: Q-Net 감정평가사 기출문제 (cst00309) artlSeq=5036790 file=extracted/1008/5036790/2011년 제22회 감정평가사 1차시험 기출문제/2교시.hwp
%% conversion: kordoc 4.24.51 -> tools/gen1.py (auto); review=proofread-gpt-astra; flags=none
\begin{question}{79}
부동산등기법령상 토지 및 건물의 등기에 관한 설명으로 옳은 것은?
\begin{choices}
  \item 토지 면적의 증감 또는 지목의 변경이 있을 때에는 그 토지 소유권의 등기명의인은 1주일 이내에 그 등기를 신청하여야 한다.
  \item 판결 또는 그 밖의 시․구․읍․면의 장의 서면에 의하여 자기의 소유권을 증명하는 자는 미등기건물의 소유권보존등기를 신청할 수 있다.
  \item 토지수용으로 인한 소유권이전등기는 토지소유자와 기업자가 공동으로 신청하여야 한다.
  \item 존재하지 아니하는 건물에 대한 등기가 있는 경우에는 등기관이 직권으로 멸실등기를 하여야 한다.
  \item 여러 개의 부동산에 관한 권리를 목적으로 하는 저당권의 설정등기를 신청하는 경우에 부동산이 3개 이상인 때에는 신청서에 공동담보목록을 첨부하여야 한다.
\end{choices}
\end{question}
\answer{2}
